\section{Catch-up, the speed limit and rents}\label{sec:catchup}

The frictionless benchmark predicts equal returns across topics. The arXiv data
reject this (Fact~\ref{fact:hump}). The same researcher earns more citations in
denser topics over most of the density range, and less in the densest ones. We
now let $b_c<\infty$: newcomers must read a topic's stock before contributing,
while its literature keeps growing.

\subsection{Mature topics}

\begin{proposition}[Mature topics under catch-up costs]\label{prop:mature}
Let \cref{ass:Q,ass:phi,ass:credit,ass:interior} hold and $b_c<\infty$.
\begin{enumerate}[label=(\roman*)]
\item $w(P)=v(P)(1-P/b_c)^\nu$ is strictly log-concave on $(0,b_c)$, with a unique peak $\hat P_w<\min\{\Phat,b_c\}$, and vanishes at $b_c$.
\item In a topic equilibrium the mature density $P^m$ solves $w(P^m)=F$ on the falling side of $w$, and $P^m<b_c$. Incumbents earn $a\,v(P^m)$ per period. Their log premium over the frontier equals the catch-up cost, $\ln[v(P^m)/F]=-\nu\ln(1-P^m/b_c)=(\rho+\xi)\,T(P^m)$.
\item Mature topics lie on the \emph{rising} side of $v$ if and only if catch-up is costly relative to the rent at the peak: $\Phat\ge b_c$, or $\nu\ln(1-\Phat/b_c)<\ln(F/\vhat)$.
\item Let topics differ in fertility $\zeta$, so that payoffs are $\zeta v(P)$, and let $P^m(\zeta)$ solve $\zeta w(P)=F$. Across mature topics, log returns rise with log density with slope equal to the catch-up elasticity,
\[
\frac{\dd\ln[\zeta\,v(P^m(\zeta))]}{\dd\ln P^m(\zeta)}=\eps_\Gamma(P^m)\equiv\nu\,\frac{P^m/b_c}{1-P^m/b_c},
\]
which exceeds the premium $-\nu\ln(1-P^m/b_c)$.
\end{enumerate}
\end{proposition}

Catch-up turns the dissipated rent into a protected one. Incumbents earn more
than the frontier, and the premium exactly offsets what a newcomer gives up
while she reads her way in. More fertile topics attract more entrants, have
larger stocks and so protect larger premia. This yields a positive
returns--density gradient across topics, whether those topics lie left or
right of $\Phat$. In the data the same gradient also reflects audience size,
since a paper with more close neighbours has more potential citers. The
measured gradient of 6.7\% per standard deviation (Fact~\ref{fact:hump}) is
therefore an upper bound on the catch-up channel. We use it as such in
\cref{sec:quant}.

\subsection{The speed limit}

\begin{theorem}[Speed limit]\label{thm:speed}
Let $b_c<\infty$.
\begin{enumerate}[label=(\roman*)]
\item \emph{(Closure.)} A topic whose flow is at least $b_c$ is closed to newcomers: anyone who joins it while its flow stays there never catches up. Its incumbents' rents are protected however large they are.
\item \emph{(Rents near the limit.)} Every mature topic satisfies $P^m<b_c$. As fertility rises, $P^m(\zeta)\to b_c$, and both the incumbents' premium and the cross-topic returns gradient diverge; the gradient grows faster than the premium.
\item \emph{(Growing topics.)} Myopic entry into a topic that grows stops before its flow reaches $b_c$, however fast it grows.
\item \emph{(Production shocks close topics.)} A production speed-up by the factor $1+\alpha$ raises the flow of every topic on impact, before membership adjusts, by the same factor at given reading shares. Every topic with $P^m>b_c/(1+\alpha)$ closes, and stays closed for at least $\ln[(1+\alpha)P^m/b_c]/\rho$ periods while turnover thins it.
\end{enumerate}
\end{theorem}

The theorem makes ``too fast'' precise. A literature can grow faster than the
people outside it can absorb, and when it does, the topic is shut to
newcomers however rewarding it is for those inside. Near the limit, small
differences in fertility produce large differences in returns, so the speed
limit is also where the returns gradient is steepest. Part (iv) is the version
that matters for artificial intelligence. A tool that lets everyone produce
faster raises every topic's flow at once, and the topics that were already
densest are pushed past the limit until turnover brings them back. Entry, and
the turnover of people and ideas that comes with it, stops there for a while.
This is one mechanism behind the slowed progress of large fields documented by
\citet{chuevans2021}: when papers arrive faster than newcomers can absorb them,
attention stays with the established canon. Better catch-up tools, such as
reviews and summaries, raise the limit.

\subsection{Rising incumbent returns, falling newcomer returns}

The interesting object is the gap between what incumbents earn, $v$, and what a
newcomer can expect over her spell, $w=v\Gamma$. Incumbents never pay the
catch-up cost again. Newcomers and the planner do. The welfare statement
follows almost by definition of stability, so we state it as a corollary and
put the weight on what it implies for measurement.

\begin{corollary}[Overcrowding with rising returns]\label{cor:overcrowd}
In every stable topic equilibrium, one more member in a topic of density $P$
changes total output of the other members by $n\,a\,w'(P)\,\dd P/\dd n<0$.
This counts both the spillover on incumbents ($v'$) and the extra catch-up she
imposes on future newcomers ($v\,\Gamma'$). If $w(\Phat)<F$, then
$\hat P_w<P^m<\Phat$: incumbents' returns rise with density ($v'(P^m)>0$) while
the topic is too crowded ($w'(P^m)<0$). The loss of output per researcher in
the large-topic limit is $1-F/\hat w$.
\end{corollary}

The measured gradient is the incumbents' private return to density. A
newcomer values a topic net of the reading she must do to join it, and so does
the planner. Rising incumbent returns are the footprint of the catch-up rent,
not of an untapped spillover. A natural reading of Fact~\ref{fact:hump} is that
most researcher-years sit where more density raises returns, so most topics are
too thin. The corollary says that this reading conflates $v$ with $w$. Two
kinds of topic can be too thin. The first is young topics that have not yet
filled. The second, if the frontier is attractive enough, is newborn topics that
never find followers (\cref{rem:coordination}).

\begin{remark}[Newborn topics]\label{rem:coordination}
A newborn topic has no stock, so its first follower needs no catch-up and
earns $v(P_0)$, where $P_0$ is the flow of a newborn topic. Newborn topics
attract followers if and only if $v(P_0)>F$. (The condition is on $v$, not on
$w$: a stationary topic of density $P_0$ would carry a stock, a newborn one
does not.) If $F>v(P_0)$, newborn topics attract no followers, although a
coordinated group would earn $\vhat>F$. Science then holds too many topics that
are too small, and the planner would fix it by seeding followers. When
$v(P_0)>F>w(P_0)$, as in the calibration, a newborn topic is worth joining
while its stock is empty, but a stationary topic of the same density is not.
The topic must grow fast enough for its stock to lag behind its flow, or
entry stalls before the topic is worth joining in its own right
(\cref{prop:toofast}). Policies that raise the frontier value by more than
$v(P_0)-F$ shut down the birth of new topics (\cref{prop:policy}).
\end{remark}

\subsection{Social value and priority transfers}

The welfare statements so far count private credit. Priority credit is partly a
transfer \citep{merton1957,hillstein2025}. When a later paper turns out to
duplicate an earlier one, some of the credit it loses goes to the earlier
paper, so part of the ``loss'' from overlap is redistribution. Conversely, a
duplicate may be worth less to science than its residual citations suggest. We
therefore separate social value from credit. Let the social value of a project
be $Q\,C_s$, with
\[
C_s(R,P)=\exp\{-\theta P[\lambda_s-\Delta_s\varphi(R/P)]-\theta_c\lambda_cP\},
\]
where $\lambda_s$ is the share of a project's social value lost per unread
covering paper and $\Delta_s$ the part of it that detection averts. Researchers
still choose $r^*(P)$ to maximise credit, and social output per unit of
production speed in a topic of density $P$ is
$s(P)=(1-r^*)\,Q(br^*,P)\,C_s(br^*,P)$, with $s_w(P)=s(P)\Gamma(P)$. The
catch-up cost is a real resource cost, reading time, and is not a transfer.

\begin{proposition}[Social value]\label{prop:social}
\begin{enumerate}[label=(\roman*)]
\item One more member changes the social output of the others by $n\,a\,s_w'(P)\,\dd P/\dd n$, where
\[
\frac{s'(P)}{s(P)}=\frac{v'(P)}{v(P)}-\theta\big[(\lambda_s-\lambda)-(\Delta_s-\Delta)\psi(c^*)\big]+\theta b(\Delta_s-\Delta)\varphi'(c^*)\,r^{*\prime}(P),
\]
with $\psi(c)=\varphi(c)-c\varphi'(c)\ge0$.
\item \emph{(Uniform wedge.)} If $\Delta_s=\Delta$, then $s(P)=v(P)e^{-\theta(\lambda_s-\lambda)P}$ exactly. Social and private values then differ by a factor that rises or falls exponentially in density. The equilibrium is overcrowded relative to the social optimum if and only if the share of overlap losses that are transfers, $\tau\equiv(\lambda-\lambda_s)/\lambda$, satisfies
\[
\tau<\frac{|\eps_w(P^m)|}{\theta\lambda P^m}.
\]
In particular overcrowding is stronger when duplicates are worth less to science than to their authors ($\lambda_s>\lambda$).
\item \emph{(Proportional transfer.)} If a share $\tau$ of every overlap loss is a transfer, $\ln C_s=(1-\tau)\ln C$, the social optimum density maximises $\ln v(P)-\tau\ln C^*(P)+\ln\Gamma(P)$, where $C^*(P)=C(br^*(P),P)$.
\end{enumerate}
\end{proposition}

Part (ii) gives a sharp test. Overcrowding survives any degree of transfer as
long as the slope of the newcomer's payoff curve at the mature density exceeds
the expected erosion intensity $\theta\lambda P^m$. In the calibrated model the
slope is \calEpsW\ against an erosion intensity of \calThetaLamPm, so the
critical transfer share is \calTauCrit, above one. Under proportional
transfers (part (iii)), the social optimum moves from \calPsocialZero\ to
\calPsocialFull\ papers a year if all overlap losses are transfers, still below
the mature density of \calPm. Transfers therefore shrink the welfare gain from
optimal topic choice, from \calPlannerGain\ to \calGainTransfer\ of output per
researcher, but they do not reverse its sign. The catch-up burden that newcomers
impose on one another is a real resource cost, and on its own it keeps mature
topics too dense. The transfer critique matters for magnitudes, and especially
for the frictionless benchmark, where the dissipated rent is a difference in
private payoffs; it does not matter for the direction of the inefficiency.
